%!TeX root=../thesis.tex

\chapter{Research Design}
\label{ch:method}
\label{chap:research-design}

This chapter describes and justifies the research design adopted in this work.
It covers the adopted research strategy along with the conceptual framework created to fulfill the research
questions. Moreover, it covers the study context and case selection, the researcher's role, the data collection procedures, and the criteria adopted to establish the trustworthiness of the findings. 
The chapter concludes with notes on the ethical considerations that governed the study.

% ---------------------------------------------------------------------------
\section{Research Strategy}
% ---------------------------------------------------------------------------

The objective of this work was to empirically evaluate and refine Hypo-Stage,
a support tool for the ArchHypo technique, by applying it with multiple software
engineering teams in a real-world organizational setting. Answering the research
questions posed in Chapter~\ref{ch:intro} requires understanding not only \emph{what}
outcomes the tool produces, but also \emph{how} teams interact with it, \emph{why}
certain difficulties arise, and \emph{which} adjustments improve its fit with existing
practices. Following the design framework for flexible qualitative inquiry \citep{robson2016}, the purpose of this work combined with its research questions clearly pointed out the need for a \textbf{flexible research design}.

Within flexible designs, the study follows a \textbf{multiple case study strategy} as defined in methodological treatments of case-study research \citep{robson2016,yin2009}.
A case study is defined by them as an empirical investigation of a particular contemporary
phenomenon within its real-life context, using multiple sources of evidence and where
the boundary between phenomenon and context is not clearly distinct. This definition
fits the present work precisely: the phenomenon under study is the use of Hypo-Stage
by a software engineering team, and its context --- the organizational, technical, and
social environment in which the team operates --- is both relevant and inseparable
from the phenomenon itself. A purely experimental design, which would require
controlled manipulation of variables and prespecification of the design before data
collection, is not appropriate here. The number of teams is small, the settings are
naturalistic and non-manipulable, and the research purpose calls for rich,
process-oriented accounts rather than statistical inference \citep{robson2016}.

The study involves \textbf{two cases}, each consisting of a distinct software
engineering team. Using multiple cases rather than a single case follows the logic
of \emph{analytic (or theoretical) generalization} from prior methodological accounts \citep{robson2016,yin2009}: the
two cases are not selected to constitute a statistical sample of teams, but to provide
complementary and contrasting instances through which patterns of tool use,
difficulties, and refinements can be examined and cross-validated. Multiple case studies are analogous to multiple
experiments where each case either replicates findings or provides theoretically predicted contrast --- both strengthening the analytic basis for generalization \citep{robson2016}.

Beyond its descriptive and explanatory purposes, the study also has an
\textbf{evaluative purpose}: it seeks to assess the worth and fitness of Hypo-Stage
as a practical tool for managing architectural uncertainty in software engineering
teams. Flexible-design accounts distinguish \emph{formative evaluation}, which
aims at understanding and improving an ongoing program or intervention, from
\emph{summative evaluation}, which renders a final verdict on overall effectiveness \citep{robson2016}.
This study is explicitly formative: data collected during the observation period are used to identify difficulties and trigger tool refinements in near real-time, rather than to produce a retrospective pass-or-fail judgment.

Finally, the iterative loop through which the researcher observes difficulties,
implements refinements to the tool, deploys an updated version, and resumes
observation incorporates a core element of \textbf{action research} as described
in the methodological literature \citep{robson2016}: a deliberate cycle of plan, act, observe,
and reflect, with the explicit intent of producing practical change alongside
research knowledge. This does not transform the study into a full action research
project --- the primary goal remains the production of knowledge about tool use and
architectural uncertainty, not organizational change more broadly --- but the
iterative refinement component follows the same logic and must be acknowledged as
part of the design.

In summary, the research strategy can be formally described as a
\textbf{flexible, multiple case study design with formative evaluation purpose and
action research elements}, following the taxonomy from \citet{robson2016}.

\subsection{Key Methodological Concepts}
\label{sec:methodological-concepts}

Because the subsequent sections, and the results chapter, rely repeatedly on a
compact set of methodological terms drawn from flexible-design guidance \citep{robson2016}, the
definitions below clarify their meaning for the remainder of this work. They are
collected here so that every later mention resolves to an unambiguous referent.

\paragraph{Flexible design:} A research design in which the detailed plan (data
sources, instruments, and categories of analysis) is allowed to evolve during
the study in response to what the researcher learns in the field, in contrast
to a \emph{fixed design}, where all elements are pre-specified before data
collection begins~\citep{robson2016}. Flexible designs are the standard choice
for case studies, ethnographies, and grounded-theory work, where the phenomenon
under study is only partially understood at the outset and the researcher
expects to refine the analytical lens as evidence accumulates.

\paragraph{Multiple case study:} An empirical investigation that examines two
or more contemporary phenomena within their real-life contexts, using multiple
sources of evidence, where the boundary between phenomenon and context is not
clearly distinct~\citep{robson2016,yin2009}. The cases are selected not as a
statistical sample but as analytical units that either replicate findings
(\emph{literal replication}) or produce theoretically predicted contrasts
(\emph{theoretical replication}).

\paragraph{Formative evaluation:} An evaluative purpose aimed at understanding
and improving an ongoing program, intervention, or artifact while it is still
being shaped, as distinguished from \emph{summative evaluation}, which delivers
a retrospective pass-or-fail verdict on overall effectiveness~\citep{robson2016}.
In the present study, the evaluative object is the Hypo-Stage tool, and the
formative stance justifies the iterative refinement cycle described in
Section~\ref{sec:procedure}.

\paragraph{Action research:} A tradition of inquiry in which the researcher
deliberately alternates between observing a setting, intervening in it, and
reflecting on the effects of that intervention, following a recurrent
\emph{plan--act--observe--reflect} cycle~\citep{robson2016}. The present study
borrows the cycle as a design element without
adopting the emancipatory or organizational-change commitments that define
full action research.

\paragraph{Purposive sampling:} A non-probabilistic sampling strategy in which
participants or cases are chosen because they are expected to be informative
about the phenomenon of interest, rather than because they are representative
of a population~\citep{robson2016}. The two teams in this study were selected
purposively to provide complementary contrast (product vs. platform scope),
not to constitute a statistical sample of engineering teams.

\paragraph{Analytic generalization:} The form of generalization proper to case
study research, in which findings are extended to theoretical propositions
rather than to populations~\citep{robson2016,yin2009}. It is contrasted with
\emph{statistical generalization}, which requires a representative sample and
probabilistic inference. The claims in Chapter~\ref{chap:results} are analytic:
they describe conditions under which the observed patterns are theoretically
expected to recur, not the frequency with which they would appear in a wider
population.

\paragraph{Member checking:} A trustworthiness procedure in which
interpretations derived by the researcher are returned to the participants
whose behavior or statements produced them, so that the participants may
verify, correct, or elaborate on the researcher's reading~\citep{robson2016}.
Section~\ref{sec:trustworthiness} describes how member checking was
operationalized in this study through the active-listening debriefs with each
team's technical leader.

These concepts recur throughout Sections~\ref{sec:research-context},
\ref{sec:data-collection}, and~\ref{sec:data-analysis}, and frame the claims
made in Chapters~\ref{chap:results} and~\ref{chap:conclusion}.

% ---------------------------------------------------------------------------
\section{Conceptual Research Framework}
% ---------------------------------------------------------------------------

The conceptual research framework organizes the key constructs of this work and their relationships in a
way that is consistent with the research questions and guides the data collection and analysis as described for flexible designs \citep{robson2016}.
Figure~\ref{fig:conceptual-framework} describes the framework.

\begin{figure}[htbp]
  \centering
  \includegraphics[width=0.95\textwidth]{conceptual-framework}
  \caption{Conceptual framework of the study.}
  \label{fig:conceptual-framework}
\end{figure}

The central piece is the \emph{Hypo-Stage Tool Usage}, understood as the practices through which
team members appropriate and practice the ArchHypo technique via Hypo-Stage: creating
and editing architectural hypotheses, assessing uncertainty levels and impact,
defining and tracking technical plans, and revisiting hypotheses over successive
sprint cycles \citep{silva2024tse}. The tool usage is conditioned by two sets of
contextual factors: the \emph{team context} (composition, experience profile, scope of
work, and familiarity with architectural concerns) and the \emph{organizational
context} (the company's development platform, the role of Backstage as the delivery
portal, and the scaled product being developed). The tool usage produces two observable
streams: the \emph{artifacts} recorded in Hypo-Stage (hypotheses, assessments,
technical plans, and their evolution over time) and the \emph{experiential observations}
of team members (benefits, difficulties, and adaptation strategies). These streams
feed the \emph{Hypo-Stage Tool Refinements} --- changes made to Hypo-Stage and its usage
guidelines --- which in turn modify subsequent tool usage, forming an iterative cycle.
This framework is consistent with the formative evaluation and action research elements of the design, and directly links each element to the research questions (RQ1--RQ4) stated in Chapter~\ref{ch:intro}.

% ---------------------------------------------------------------------------
\section{Research Context}
\label{sec:research-context}
% ---------------------------------------------------------------------------

The study was conducted within a technology organization that operates
large-scale software products across Latin America and employs over 500 software
engineers. The organization is an intensive user of Backstage\footnote{\url{https://github.com/backstage/backstage}}, an open-source
developer portal, which serves as the central platform for internal tooling,
documentation, and engineering workflows. The availability of this platform was an
enabling condition for the study: Hypo-Stage was deployed as a Backstage plugin, so its adoption is conditioned by the availability of that platform.
The Hypo-Stage tool was deployed directly in the company's production environment, integrated with the same catalog and workflow interface that engineers use daily, reducing adoption friction considerably. 
This allowed the researcher to observe the tool usage in a naturalistic setting, with minimal interference from the researcher.

\subsection{Case Selection}
\label{sec:case-selection}

Two teams were selected as cases. Following methodological guidance on purposive sampling in flexible designs \citep{robson2016}, the cases were not chosen
to be statistically representative of a population of teams, but to provide
\emph{complementary contrast}: teams with different compositions, scopes, and
relationships to the broader system, so that differences and similarities in tool use
can be analytically compared.

\begin{description}
  \item[Team~1 --- Product Team.] A team of 12 software engineers (5 senior,
    7 mid-level) responsible for developing new product features with direct
    impact on the company's business objectives and end-user experience. This team
    works on a bounded product domain, with close proximity to business requirements
    and user feedback. Its architectural concerns are primarily about supporting new
    features reliably and at scale while preserving existing quality attributes.

  \item[Team~2 --- Platform Team.] A team of 4 software engineers (2 senior,
    2 mid-level) responsible for developing and maintaining broad-scope internal
    tooling that supports software development across the entire organization. This
    team's architectural decisions affect how all engineers in the company build,
    deploy, and operate software. Its architectural concerns are therefore more
    horizontal, with higher impact on organizational engineering practices and
    greater uncertainty about usage patterns and requirements.
\end{description}

\paragraph{Literal and theoretical replication.} In the multiple-case replication logic for case studies \citep{yin2009}, two or more cases are deliberately selected to perform one of
two complementary analytical roles. \emph{Literal replication} chooses cases
that share the relevant conditions and are therefore predicted to produce
similar findings; convergence across literal replications strengthens the
claim that the finding is genuinely tied to the intervention rather than to
idiosyncratic case features. \emph{Theoretical replication} chooses cases
that differ on a theoretically meaningful dimension and are therefore
predicted to produce contrasting findings; divergence along the predicted
direction strengthens the theoretical account of \emph{why} the intervention
behaves differently across conditions.

The two cases represent a \emph{literal replication} in so far as both teams face
the same intervention (use of Hypo-Stage and ArchHypo) and a \emph{theoretical
replication} in that their different scopes and compositions should, theoretically,
lead to different experiences of architectural uncertainty and different challenges
in adopting a hypothesis-driven approach \citep{robson2016,yin2009}. This contrast
strengthens the analytic basis for identifying which findings are context-independent and which are conditioned by team type.

\begin{table}[htbp]
\centering
\caption{Case comparison: Team 1 (Product) and Team 2 (Platform).}
\label{tab:case-comparison}
\begin{tabular}{p{0.30\linewidth} p{0.32\linewidth} p{0.32\linewidth}}
\toprule
\textbf{Dimension} & \textbf{Team 1 --- Product Team} & \textbf{Team 2 --- Platform Team} \\
\midrule
Size & 12 engineers (5 senior, 7 mid-level). & 4 engineers (2 senior, 2 mid-level). \\
\midrule
Scope & Bounded product domain; new features with direct business and user impact. & Broad-scope internal tooling supporting all engineers in the organization. \\
\midrule
Architectural concerns & Supporting new features reliably at scale; preserving existing quality attributes. & Horizontal decisions affecting build, deploy, and operate practices organization-wide. \\
\midrule
Uncertainty profile (expected) & Moderate uncertainty, high proximity to user feedback. & Higher uncertainty about usage patterns; high cross-team impact of decisions. \\
\midrule
Replication logic & Literal replication --- same intervention (Hypo-Stage + ArchHypo). & Theoretical replication --- different scope and composition expected to surface different challenges. \\
\bottomrule
\end{tabular}
\end{table}

% ---------------------------------------------------------------------------
\section{Researcher's Role}
\label{sec:researcher-role}
% ---------------------------------------------------------------------------

Throughout the study, the researcher occupied the dual role of
\textbf{participant observer and facilitator}, in the terms used in flexible-design methodology \citep{robson2016}.
The researcher is an employee of the organization in which the study was conducted, which makes this an instance of \emph{insider research}.
This position provided privileged access to the teams, their work context, and their daily practices, which would have been substantially harder to
obtain as an external investigator. It also allowed the researcher to introduce
Hypo-Stage into the company's infrastructure in a credible and practically feasible
way.

As discussed in accounts of insider research \citep{robson2016}, such studies raise specific concerns
regarding bias and the potential conflation of the researcher's own perspective with
those of participants. These concerns are addressed in
Section~\ref{sec:trustworthiness}. It is important to note that, while the researcher
served as facilitator during training sessions and provided assistance on request
throughout the observation period, he did not direct the teams in \emph{which}
architectural uncertainties to address or \emph{how} to formulate specific
hypotheses. The content of the work produced by the teams in Hypo-Stage reflects
their own judgments and priorities.

% ---------------------------------------------------------------------------
\section{Procedure}
\label{sec:procedure}
% ---------------------------------------------------------------------------

\subsection{Training and Introduction}

Before the observation period began, each team participated in a one-hour training
session led by the researcher. The session covered:

\begin{itemize}
  \item[(i)] The conceptual foundations of hypotheses engineering and the ArchHypo technique, including the distinction between architectural hypotheses, uncertainty levels, impact assessment, and technical plans.
  \item[(ii)] The context and motivation of the research, including the broader challenge of making architectural uncertainties explicit in teams working on scaled products.
  \item[(iii)] A live demonstration of Hypo-Stage running in the company's Backstage instance, illustrating how to create and manage hypotheses within the tool. Teams were given written communication prior to the meeting (Appendix~\ref{apx:training-communication}) --- including an illustration of the ArchHypo big picture (Figure~\ref{fig:archhypo-overview}) --- so that they could arrive at the session with some prior exposure to the key concepts.
\end{itemize}

\begin{figure}[htbp]
  \centering
  \includegraphics[width=0.85\textwidth]{archhypo-big-picture}
  \caption{Overview of the ArchHypo technique showing the iterative cycle through
    which engineers identify architectural hypotheses, assess their uncertainty
    and impact, define technical plans, and take or postpone architectural
    decisions. Source: adapted from \citet{silva2024tse}.}
  \label{fig:archhypo-overview}
\end{figure}

\subsection{Observation Period}

Following the training session, each team entered an observation period during which they were invited to use Hypo-Stage to
register and manage any architectural uncertainties that arose in the course of their
regular development work. The organization operates on two-week agile sprints. The initial participant communication (Appendix~\ref{apx:training-communication}) described a four-week horizon; in practice, teams continued to use the tool voluntarily beyond that window.

The empirical window for hypothesis creation is anchored in Hypo-Stage registration timestamps. As documented in Appendix~\ref{app:hypotheses-data}, the thirteen hypotheses were created between 11~February~2026 and 12~April~2026: an interval of 61~calendar days (approximately nine weeks), spanning five consecutive two-week sprint periods from the earliest to the latest registration. No constraints were imposed on the type, number, or
formulation of hypotheses, nor on when to create them during the sprints. It is intrisic to the nature of the organization in which the experience was conducted that the teams are free to adjust their agile practices to their own needs. Therefore, in contrast to the PropSEP evaluation setting \citep{silva2024tse}, the teams were not required to create hypotheses during the sprints, but rather to use the tools whenever they deemed necessary.
The objective was to observe tool use in naturalistic conditions, with minimal researcher interference in the content of the teams' work.

During this period, the researcher remained available to both teams for support on
request. Support was of two distinct types:

\begin{enumerate}
  \item \textbf{Conceptual support.} When a team member found it difficult to
    formulate a hypothesis or to assess its uncertainty and impact within the
    ArchHypo framework, the researcher provided clarification or facilitated a brief
    discussion. Interactions of this type were noted as observations about the
    conceptual difficulty of the technique.

  \item \textbf{Tool support.} When a team member encountered a difficulty that
    originated from the tool itself --- a usability problem, a missing feature, or a
    mismatch between the tool's behavior and the expected workflow --- the researcher
    logged this as a tool improvement candidate and, as a matter of priority,
    implemented and deployed an updated version of Hypo-Stage as quickly as
    practicable. This was done to ensure that tool limitations would not persistently
    obstruct teams' ability to practice ArchHypo, and to allow subsequent
    observations to reflect a more refined version of the tool.
\end{enumerate}

This second type of support operationalizes the \emph{action research cycle} within
the study: difficulty observed $\rightarrow$ refinement designed $\rightarrow$
updated tool deployed $\rightarrow$ use and observation resumed. Each refinement
cycle is documented in Section~\ref{sec:refinements}.

% ---------------------------------------------------------------------------
\section{Data Collection}
\label{sec:data-collection}
% ---------------------------------------------------------------------------

\begin{table}[htbp]
\centering
\caption{Data collection instruments, sources, and research questions addressed.}
\label{tab:data-collection}
\begin{tabular}{p{0.25\linewidth} p{0.28\linewidth} p{0.20\linewidth} p{0.19\linewidth}}
\toprule
\textbf{Instrument} & \textbf{Data Produced} & \textbf{Collection Timing} & \textbf{RQs Addressed} \\
\midrule
Participant observation (field notes) & Reactions, use patterns, informal commentary, support requests. & Continuous during observation period. & RQ1, RQ3 \\
\midrule
Artifact analysis (Hypo-Stage records) & Hypotheses, assessments, technical plans, lifecycle transitions. & End of observation period. & RQ1, RQ2 \\
\midrule
Open-ended active listening sessions & Perceived usefulness, difficulties, suggestions, overall assessment. & End of observation period. & RQ3, RQ4 \\
\bottomrule
\end{tabular}
\end{table}

Following the multiple data source approach recommended for case studies \citep{robson2016,yin2009}, data were collected through three complementary channels:

\subsection{Participant Observation}
\label{subsec:participant-observation}

The researcher's continuous presence in the organization and involvement in both
teams' environments provided an ongoing stream of observational data. Observations
covered: 

\begin{itemize}
  \item[(i)] Team members' reactions and difficulties during and after the training
  session;
  \item[(ii)] Patterns of use (or non-use) of Hypo-Stage during the sprint;
  \item[(iii)] Interactions in which team members requested conceptual or tool support;
  \item[(iv)] Informal commentary about the fit of ArchHypo with existing practices.
\end{itemize}

Following established guidance on participant observation in flexible designs \citep{robson2016}, field notes were recorded as close as possible in time to the observed events, and the researcher maintained a reflexive awareness of his dual role as insider and investigator.

\subsection{Artifact Analysis}
\label{subsec:artifact-analysis}

The hypotheses, assessments, and technical plans created by each team in Hypo-Stage
constitute a documentary record of how they used the tool over the observation period.
These artifacts were collected at the end of the observation period and analyzed to
characterize:

\begin{itemize}
  \item[(i)] The number and types of hypotheses created;
  \item[(ii)] The quality attributes and sources of uncertainty they addressed;
  \item[(iii)] How hypotheses were assessed and what technical plan actions were defined;
  \item[(iv)] How hypotheses evolved (created, refined, or closed) across the observation period (five two-week sprints spanned by registration timestamps; Appendix~\ref{app:hypotheses-data}).
\end{itemize}

The content analysis of documentary record is a recognized data collection method in flexible case study designs \citep{robson2016}.

\subsection{Open-Ended Active Listening Sessions}
\label{subsec:active-listening}

In addition to informal interactions arising from support requests, the researcher
conducted brief open-ended active listening sessions with team members at the end of the
observation period, inviting them to reflect on their experience of using Hypo-Stage
and ArchHypo. These conversations followed a semi-structured format, covering the
perceived usefulness of the technique, the difficulties encountered, the changes they
would suggest to the tool or its integration into their workflow, and their overall
assessment of the value of making architectural uncertainties explicit in this way.
Semi-structured or focused interviews are a central method for flexible design
research and are particularly appropriate when the researcher seeks to capture
participants' perspectives and interpretations rather than confirming predetermined
hypotheses as described for flexible qualitative designs \citep{robson2016}.

% ---------------------------------------------------------------------------
\section{Data Analysis}
\label{sec:data-analysis}
% ---------------------------------------------------------------------------

Data from the three collection sources described in
Section~\ref{sec:data-collection} were analyzed using \emph{thematic coding
analysis} as specified for qualitative data gathered under a flexible design \citep{robson2016}. Thematic coding is a family of qualitative analytic
techniques that produces a structured account of what is present textual or observational data by iteratively assigning short labels ---
\emph{codes} --- to meaningful segments, aggregating those codes into
higher-order \emph{categories} or \emph{themes}, and then comparing the
resulting structure across units of analysis~\citep{robson2016}. It is
distinguished from fully inductive \emph{grounded theory} by its willingness
to accept a priori analytical constructs (in this study, those of the
conceptual framework in Figure~\ref{fig:conceptual-framework}) and from fully
deductive \emph{content analysis} by its commitment to letting the codebook
emerge from the data rather than being fixed in advance.

In line with established guidance on flexible designs \citep{robson2016}, the analysis proceeded
through three stages, each answering a distinct analytical question.

\paragraph{Stage 1 --- Inductive open coding:} All field notes, extracted
Hypo-Stage artifact records, and active-listening transcripts were reviewed
line-by-line. A short descriptive \emph{code} was assigned to each segment that carried
analytical weight, capturing a recurring observation, a reported experience,
or a notable silence. Codes at this stage are descriptive, not interpretive:
the aim is faithful, low-inference labeling of what is present in the data.
No a priori categories were imposed. Representative codes and their
traceability to later sections appear in Table~\ref{tab:code-traceability}.

\paragraph{Stage 2 --- Axial grouping into conceptual categories:} The
descriptive codes produced in Stage~1 were grouped into higher-order
\emph{axial categories} that correspond to the constructs of the conceptual
framework. The term \emph{axial} follows the Strauss--Corbin tradition adopted
in accounts of thematic coding under flexible designs \citep{robson2016}: it denotes the relational re-organization of open codes
around a central analytical axis, in this study the constructs of tool usage,
architectural uncertainty, perceived benefits, perceived challenges, candidate
refinements, and organizational-adoption conditions. Six axial categories
emerged, each corresponding to one construct in
Figure~\ref{fig:conceptual-framework} or to an axis along which the
framework's observables were stratified.

\paragraph{Stage 3 --- Cross-case comparison:} The axial categories were then
examined across the two cases (Product Team and Platform Team) to identify
findings common to both teams and findings that were specific to one.
Similarities are reported as shared observations in
Chapter~\ref{chap:results}; differences are attributed to the distinct scopes
and compositions of product versus platform work, following the logic of
literal and theoretical replication set out in
Section~\ref{sec:case-selection}.

The complete mapping from representative open codes to axial categories,
research questions, and chapter sections appears in
Table~\ref{tab:code-traceability} at the start of
Chapter~\ref{chap:results}. That table is the primary audit trail linking
raw observations to every empirical claim made in this thesis.

Consistent with the observation that flexible designs may
include modest quantitative summaries without becoming multi-strategy
designs \citep{robson2016}, the artifact data extracted from Hypo-Stage were also summarized
descriptively (counts of hypotheses per team, frequency of quality
attributes, technical-plan strategy distribution). These summaries
complement and illustrate the qualitative findings; they do not constitute
a separate quantitative strand of analysis.

% ---------------------------------------------------------------------------
\section{Establishing Trustworthiness}
\label{sec:trustworthiness}
% ---------------------------------------------------------------------------

\begin{table}[htbp]
\centering
\caption{Trustworthiness criteria and strategies applied in this study \citep{robson2016}.}
\label{tab:trustworthiness}
\begin{tabular}{p{0.20\linewidth} p{0.35\linewidth} p{0.37\linewidth}}
\toprule
\textbf{Criterion} & \textbf{Meaning} & \textbf{Strategy Applied} \\
\midrule
Credibility & Internal validity; do findings accurately represent participants' realities? & Prolonged engagement ($\approx$9 weeks, five two-week sprints covered by registrations); triangulation across three data sources; member checking. \\
\midrule
Transferability & Degree to which findings may inform other settings. & Thick description of context; analytic (not statistical) generalization claimed. \\
\midrule
Dependability & Reliability and consistency of the research process. & Audit trail: dated field notes, systematic artifact extraction, commits linked to each refinement. \\
\midrule
Confirmability & Degree to which findings reflect data rather than researcher bias. & Explicit role separation; active search for disconfirming evidence; all claims grounded in specific observations or artifacts. \\
\bottomrule
\end{tabular}
\end{table}

The trustworthiness of findings from flexible design research is discussed in terms of \emph{credibility}, \emph{transferability},
\emph{dependability}, and \emph{confirmability} \citep{robson2016}. Each is addressed
below.

\textbf{Credibility} --- the internal validity of the findings --- was supported
through: 
\begin{itemize}
  \item[(i)] Prolonged engagement with both teams over approximately nine calendar weeks of real work in their natural environment (Hypo-Stage registrations span five consecutive two-week sprints; Appendix~\ref{app:hypotheses-data});
  \item[(ii)] Multiple data sources (observation, artifact analysis, active listening sessions), providing triangulation of evidence;
  \item[(iii)] \textbf{Member checking.} Interpretations that the researcher
    formed during the observation period --- for instance, the inference that
    hypothesis registration clustered around protected discussion time ---
    were surfaced back to each team's senior technical leader during the
    active-listening sessions described in
    Section~\ref{subsec:active-listening}. Each leader was explicitly invited
    to confirm, correct, or elaborate on the researcher's reading. This
    operationalizes the member-checking procedure recommended as a core credibility control for flexible qualitative
    designs \citep{robson2016}.
\end{itemize}

\textbf{Transferability} --- the degree to which the findings may inform other
settings --- is constrained by the specifics of the organizational context (a large
technology company in Latin America with an established Backstage platform) and the
tool (Hypo-Stage). The study does not claim statistical generalization. Rather, it
offers analytic generalization \citep{robson2016,yin2009}: the patterns of use,
difficulties, and refinements documented here provide a theoretically informed
account that researchers and practitioners can assess for relevance to their own
contexts.

\textbf{Dependability} --- the reliability of the research process --- was addressed
by maintaining an audit trail: field notes were dated and linked to specific events,
artifact data were extracted and stored systematically, and each tool refinement was
documented with its triggering observation and the corresponding commit in the
Hypo-Stage repository.

\textbf{Confirmability} --- the degree to which findings reflect the data rather
than researcher bias --- was addressed by: 

\begin{itemize}
  \item[(i)] Explicitly distinguishing between the researcher's facilitation role (which involved helping teams understand the technique) and his investigator role (which required non-directive observation of how they used it);
  \item[(ii)] Seeking disconfirming evidence, that is, noting instances where teams found value in the tool or the technique as readily as instances of difficulty;
  \item[(iii)] Grounding all claims in specific observations, artifacts, or conversation excerpts, reported in Chapter~\ref{chap:results} (including Section~\ref{sec:refinements} where refinements are documented).
\end{itemize}

% ---------------------------------------------------------------------------
\section{Ethical Considerations}
\label{sec:ethical-considerations}
% ---------------------------------------------------------------------------

The study raises ethical considerations associated with insider research \citep{robson2016}. Participation was voluntary: team members were informed of the research objectives, were free to use Hypo-Stage as much or as little as they chose, 
and were invited but not required to engage in active listening sessions. No personal performance data were collected or reported; all analysis
pertains to team-level patterns and tool-level observations. The organization in which the study was conducted and the teams involved are anonymized in reporting.
As no sensitive personal data were collected, no ethical review was required for this study.

The dual role of the researcher as organizational employee and academic investigator is a potential source of role conflict. This was managed by maintaining a clear
separation between the researcher's responsibilities to the organization (deploying and supporting a useful tool) and his responsibilities as a researcher (reporting
findings honestly, including negative ones). The research was conducted with the knowledge and informal consent of the relevant engineering leadership in the organization.
